\documentclass[11pt]{article}
\usepackage[margin=1in]{geometry}
\usepackage{booktabs}
\usepackage{url}
\title{Why the noisy 15-m/s braking-lead encounter collides}
\date{Prepared October 4, 2026}
\begin{document}
\maketitle

\textbf{Case.} The case is the controller at abd8740 (direct issue, plan-innovation trust) with
the noisy estimator (\texttt{estimatorControllerIntegrationConfig}, seed
20261003), 15~m/s brakingLead. The ego passes a lead that brakes at 0.5~m/s$^2$. The rectangles touch at 3.087~s,
hold 62. The run reproduces exactly with per-hold plan logging. No production
source was changed; the instrumentation is in \path{instrumentation.diff}.

\section{What is not the cause}
\textbf{Estimation.} The ego position error is at most 0.01~m and its yaw error
at most 0.071~rad. The target position error reaches 0.47~m. Nevertheless, the
estimated gap tracks the true gap within 0.07~m (0.044 versus 0.050~m at
3.05~s): the controller knew the gap was closing.

\textbf{Observer tightening.} The ego posterior generator is zero and the
collision tightening is 0 throughout. The observer's uncertainty did not enter
this problem.

\textbf{Trust scale as trigger.} At the trigger (2.30~s) the trust scale was 1.

\section{Causal chain}
\begin{enumerate}
\item From 1.90 to 2.25~s, the shifted overtaking plan holds a predicted gap of
  exactly 0.200~m. Its endpoint deviation from the terminal base grows: in
  scaled units it is 0.88 at 2.25~s and 0.96 at 2.30~s, against a terminal core
  radius of 0.001.
\item At 2.30~s the shifted problem fails in both stages (CLF without a
  numerical result, then PCBF restoration infeasible). A replay of that hold
  shows the cause. Removing only the terminal cone makes the shifted plan
  solvable; it issues steering $-0.142$~rad and almost no braking, so it
  continues the overtake. Removing the collision rows, path rows, terminal
  target rows or state limits does not help.
\item The fallback potential-field seed has an exact terminal endpoint. Its
  plan brakes (ratio 0.149) and later steers back toward the lane. It is
  feasible only with positive prefix safety slack.
\item The braking transition produced a 0.22~m model innovation, and the trust
  scale fell to 0.08--0.15. Every later hold repeated the same chain:
  terminal-infeasible shift, then a fresh seed with small trust. From 2.70~s
  each issued plan predicted contact itself (predicted minimum gap 0, slack
  about 3~m). Contact occurred at 3.09~s.
\end{enumerate}

\section{Counterfactual closed loops}
\begin{center}
\begin{tabular}{lll}
\toprule
Variant & outcome & minimum gap\\
\midrule
adaptive trust (as committed) & collision at hold 62 & 0\\
trust fixed 1 & no solution at hold 19 & 1.402~m\\
trust floor 0.5 & no solution at hold 24 & 1.390~m\\
trust floor 0.25 & no solution at hold 35 & 1.358~m\\
adaptive trust, terminal cone removed & no collision, unrecovered at 100~s & 0.068~m\\
trust fixed 1, terminal cone removed & recovered at hold 533 & 0.040~m\\
\bottomrule
\end{tabular}
\end{center}

All variants with the terminal cone fail. A trust floor only changes the
failure from a collision to an earlier stop. Without the terminal cone neither
run collides, although both pass within 4--7~cm, below the 0.20~m planning
margin.

\section{Interpretation}
The primary trigger is the terminal requirement. The endpoint of an ongoing
overtake must coincide with the cruise equilibrium to within a core radius of
0.001. Once the shifted plan cannot meet that within its input box, the
controller abandons the overtake for a potential-field seed that returns to
the lane. Collision rows on the prefix are soft, so the solver then accepts
predicted contact instead of reporting failure. The small adaptive trust
scale keeps later plans close to those seeds; with a larger floor, the same
chain ends in an infeasible hold instead. The same collision-row/terminal-cone
conflict stopped the exact-observation 8-m/s brakingLead at hold 108.

No change is proposed or validated here. Removing the terminal cone affects
the recursive-feasibility argument and was tested only as a counterfactual.
The files in \path{COLLISION_DIAGNOSIS_20261004/} are the variant outcomes and
the run, analysis and replay scripts (stored as \path{.m.txt}).
\end{document}
